\section{Preparations using free probability}

Let $(M,\tau)$ be a finite von Neumann algebra and we will typically assume that it is a ${\rm II}_1$-factor. The $2$-norm on $M$ is defined as $\|x\|_2 = \tau(x^*x)^{1/2}$. We denote its unitary group by ${\rm U}(M)$ and its projective unitary group by ${\rm PU}(M):={\rm U}(M)/S^1$. Consider unitaries $u,v \in {\rm U}(M)$ which are freely independent, that is to say that they lie in freely independent subalgebras in the sense of Voiculescu \cite{MR1217253}. We set $\alpha := \tau(u)$ and $\beta := \tau(v)$. Note that because of freeness we have
\begin{gather*}
0 = \tau((u-\alpha)(v-\beta)) = \tau(uv) - \beta \tau(u) - \alpha \tau(v) + \alpha \beta
\end{gather*}
and hence $\tau(uv) = \tau(u)\tau(v)$ for any freely independent pair of unitaries $u,v \in {\rm U}(M)$. We aim to compute the trace of $uvu^*v^*$. Again, by freeness, we obtain
\begin{align*}
0 ={}& \tau((u-\alpha)(v-\beta)(u^* - \bar \alpha)(v^* - \bar \beta))
\\
={}& \tau(uvu^*v^*) - \alpha \tau(vu^*v^*) - \beta \tau(uu^*v^*) - \bar \alpha \tau(uvv^*) - \bar \beta \tau(uvu^*)
\\
& + \alpha \beta \tau(u^*v^*) + |\alpha|^2 \tau(vv^*) + \alpha \bar \beta \tau(vu^*) + \bar \alpha \beta \tau(uv^*) + |\beta|^2 \tau(uu^*) + \bar \alpha \bar \beta \tau(uv)
\\
& - |\alpha|^2 \beta \tau(v^*) - \alpha |\beta|^2 \tau(u^*) - |\alpha|^2 \bar \beta \tau(v) - \bar \alpha |\beta|^2 \tau(u) + |\alpha|^2|\beta|^2
\\
={}& \tau(uvu^*v^*) - 2|\alpha|^2 - 2 |\beta|^2 + 4 |\alpha|^2 |\beta|^2 + |\alpha|^2 + |\beta|^2 - 4 |\alpha|^2 |\beta|^2 + |\alpha|^2 |\beta|^2
\\
={}& \tau(uvu^*v^*) - |\alpha|^2 - |\beta|^2 + |\alpha|^2 |\beta|^2.
\end{align*}
Hence, we conclude that
\[
\tau(uvu^*v^*) = |\alpha|^2 + |\beta|^2 - |\alpha|^2 |\beta|^2 = 1 - \big(1-|\tau(u)|^2\big)\big(1-|\tau(v)|^2\big)
\] for any freely independent pair of unitaries $u,v \in {\rm U}(M)$.
On ${\rm U}(M)$, we consider the length function $\ell(u):= \|1-u\|_2 = \sqrt{2-2 \operatorname{Re}(\tau(u))}$.
Since
\begin{gather*}
(1-|\tau(u)|^2) = (1+ |\tau(u)|)(1-|\tau(u)|) \leq 2 (1-|\tau(u)|) \leq 2 - 2\operatorname{Re}(\tau(u)) = \ell(u)^2,
\end{gather*}
we obtain
\begin{align*}
\ell([u,v]) &= \sqrt{2-2 \operatorname{Re}(\tau(uvu^*v^*))} \\
&= \sqrt{2\big(1-|\tau(u)|^2\big)\big(1- |\tau(v)|^2\big)} \\
&\leq\sqrt{2} \cdot \ell(u) \ell(v)
\end{align*}
for any freely independent pair of unitaries $u,v \in {\rm U}(M)$. With only little more work, we also obtain a lower bound for $\ell([u,v])$.
Indeed, define the projective length function
\begin{gather*}
\bar \ell(u) = \inf_{|\lambda|=1 } \|\lambda - u\|_2 = \sqrt{2(1-|\tau(u)|)},
\end{gather*}
which metrizes ${\rm PU}(M)$ in a natural way. Note that
$\mu - \mu^2/4 \geq 1/2 \cdot \mu$ for $\mu \in [0,2]$. If $u$, $v$ are freely independent unitaries, we get
\begin{align*}
\ell([u,v])^2 &= 2(1 - \operatorname{Re} \tau([u,v])) \\
&= 2 \big(1-|\tau(u)|^2\big)\big(1- |\tau(v)|^2\big) \\
&= 2 \big(1-\big(1-\bar\ell(u)^2/2\big)^2\big) \big(1-\big(1-\bar\ell(v)^2/2\big)^2\big) \\
&= 2 \big( \bar \ell(u)^2 - \bar \ell(u)^4/4\big) \big( \bar \ell(v)^2 - \bar \ell(v)^4/4\big) \\
&\geq 2 (1/2)^2 \cdot \ell(u)^2 \ell(v)^2.
\end{align*}
Hence, we conclude that
\begin{gather*}
\ell([u,v]) \geq 1/\sqrt{2} \cdot \bar \ell(u) \bar \ell(v)
\end{gather*}
and we may summarize the computations as follows:
\begin{Theorem}
Let $u,v \in {\rm U}(M)$ be freely independent unitaries. Then, we have
\begin{gather*}
1/\sqrt{2} \cdot \bar \ell(u) \bar \ell(v) \leq \bar\ell([u,v])= \ell([u,v]) \leq \sqrt{2} \cdot \bar \ell(u) \bar \ell(v) \leq \sqrt{2} \cdot \ell(u) \ell(v).
\end{gather*}
\end{Theorem}

For the rest of the paper we fix a sequence of words $(w_n)_n$ in $F_2= \langle x,y \rangle$ such that $w_1=x$ and inductively $w_{n+1} = [w_n,y^{n}xy^{-n}]$, i.e., $w_2=\big[x,yxy^{-1}\big]$, $w_3=\big[\big[x,yxy^{-1}\big],y^2xy^{-2}\big],$ and so~on. If $u$, $v$ are freely independent and $v$ is a Haar unitary (i.e., a unitary with $\tau(v^n)=0$ for all $n \in \mathbb Z \setminus \{0\}$), it is easy to see that $ u,vuv^*,v^2u(v^*)^2,\dots$ are freely independent too. We consider the sequence $(w_n(u,v))_n$, for $i \in \mathbb N$, i.e.,
\begin{gather*}
u_1=u, \qquad
u_2=[u,vuv^*], \qquad
u_3=\big[[u,vuv^*],v^2u(v^*)^2 \big],\qquad \dots\,.
\end{gather*}
Thus, we obtain from the previous theorem by induction
\begin{gather*}
\big(1/\sqrt{2}\big)^{n-1} \cdot \bar\ell(u)^n \leq \ell(w_n(u,v)) \leq \big(\sqrt{2}\big)^{n-1} \cdot \ell(u)^n.
\end{gather*}

The upper bound was rather surprising for us at first, since $\ell(u)< 1/\sqrt{2}$ does not seem to be a severe assumption and we may assume in addition that the subgroup generated by $u$ in ${\rm PU}(M)$ is discrete, we may even assume that $u'$ is a Haar unitary in a corner $pMp$ and $u = u' + p^{\perp}$ for a projection $p$ with $\tau(p)<1/4$.

